% chapter05.tex — Limits
\spivakchapter{5}{Limits}

The concept of a limit is surely the most important, and probably the most
difficult one in all of calculus. The goal of this chapter is the definition of
limits, but we are, once more, going to begin with a provisional definition;
what we shall define is not the word ``limit'' but the notion of a function
approaching a limit.

\begin{spivakdefn}[PROVISIONAL DEFINITION]
The function $f$ approaches the limit $l$ near $a$, if we can make $f(x)$ as
close as we like to $l$ by requiring that $x$ be sufficiently close to, but
unequal to, $a$.
\end{spivakdefn}

Of the six functions graphed in Figure~1, only the first three approach $l$ at
$a$. Notice that although $g(a)$ is not defined, and $h(a)$ is defined ``the
wrong way,'' it is still true that $g$ and $h$ approach $l$ near $a$. This is
because we explicitly ruled out, in our definition, the necessity of ever
considering the value of the function at $a$---it is only necessary that $f(x)$
should be close to $l$ for $x$ close to $a$, but \emph{unequal to} $a$. We are
simply not interested in the value of $f(a)$, or even in the question of
whether $f(a)$ is defined.

\begin{center}
\figtex{assets/geometric/ch05-fig01}\\[0.6em]
\figlabel{1}
\end{center}
One convenient way of picturing the assertion that $f$ approaches $l$ near $a$
is provided by a method of drawing functions that was not mentioned in
Chapter~4. In this method, we draw two straight lines, each representing
$\mathbf{R}$, and arrows from a point $x$ in one, to $f(x)$ in the other.
Figure~2 illustrates such a picture for two different functions.

\begin{center}
\figtex{assets/geometric/ch05-fig02}\\[0.6em]
\figlabel{2}
\end{center}
Now consider a function $f$ whose drawing looks like Figure~3. Suppose we ask
that $f(x)$ be close to $l$, say within the open interval $B$ which has been
drawn in Figure~3. This can be guaranteed if we consider only the numbers $x$
in the interval $A$ of Figure~3. (In this diagram we have chosen the largest
interval which will work; any smaller interval containing $a$ could have been
chosen instead.) If we

\begin{center}
\figtex{assets/geometric/ch05-fig03}\\[0.6em]
\figlabel{3}
\end{center}
\begin{center}
\figtex{assets/geometric/ch05-fig04}\\[0.6em]
\figlabel{4}
\end{center}
choose a smaller interval $B'$ (Figure~4) we will, usually, have to choose a
smaller $A'$, but no matter how small we choose the open interval $B$, there is
always supposed to be some open interval $A$ which works.

A similar pictorial interpretation is possible in terms of the graph of $f$,
but in this case the interval $B$ must be drawn on the vertical axis, and the
set $A$ on the horizontal axis. The fact that $f(x)$ is in $B$ when $x$ is in
$A$ means that the part of the graph lying over $A$ is contained in the region
which is bounded by the horizontal lines through the end points of $B$;
compare Figure~5(a), where a valid interval $A$ has been chosen, with
Figure~5(b), where $A$ is too large.

\begin{center}
\figtex{assets/geometric/ch05-fig05}\\[0.6em]
\figlabel{5}
\end{center}
To take a specific simple example, let's consider the function $f(x)=3x$ with
$a=5$ (Figure~6). Presumably $f$ should approach the limit $15$ near $5$---we
ought to be able to get $f(x)$ as close to $15$ as we like if we require that
$x$ be sufficiently close to $5$. To be specific, suppose we want to make sure
that $3x$ is within $\frac{1}{10}$ of $15$. This means that we want to have
\[
15 - \frac{1}{10} < 3x < 15 + \frac{1}{10},
\]
which we can also write as
\[
-\frac{1}{10} < 3x - 15 < \frac{1}{10}.
\]
To do this we just have to require that
\[
-\frac{1}{30} < x - 5 < \frac{1}{30},
\]
or simply $\abs{x-5}<\frac{1}{30}$; There is nothing special about the number
$\frac{1}{10}$. It is just as easy to guarantee that
$\abs{3x-15}<\frac{1}{100}$; simply require that
$\abs{x-5}<\frac{1}{300}$. In fact, if we take any positive number
$\varepsilon$ we can make $\abs{3x-15}<\varepsilon$ simply by requiring that
$\abs{x-5}<\varepsilon/3$.

\begin{center}
\figtex{assets/plots/ch05-fig06}\\[0.6em]
\figlabel{6}
\end{center}
There's also nothing special about the choice $a=5$. It's just as easy to see
that $f$ approaches the limit $3a$ at $a$ for any $a$: To ensure that
\[
\abs{3x-3a}<\varepsilon
\]
we just have to require that
\[
\abs{x-a}<\frac{\varepsilon}{3}.
\]

Naturally, the same sort of argument works for the function
$f(x)=3{,}000{,}000x$. We just have to be $1{,}000{,}000$ times as careful,
choosing $\abs{x-a}<\varepsilon/3{,}000{,}000$ in order to ensure that
$\abs{f(x)-a}<\varepsilon$.

The function $f(x)=x^2$ is a little more interesting. Presumably, we should be
able to show that $f(x)$ approaches $9$ near $3$. This means that we need to
show how to ensure the inequality
\[
\abs{x^2-9}<\varepsilon
\]
for any given positive number $\varepsilon$ by requiring $\abs{x-3}$ to be
small enough. The obvious first step is to write
\[
\abs{x^2-9}=\abs{x-3}\cdot\abs{x+3},
\]
which gives us the useful $\abs{x-3}$ factor. Unlike the situation with the
previous examples, however, the extra factor here is $\abs{x+3}$, which isn't
a convenient constant like $3$ or $3{,}000{,}000$. But the only crucial thing
is to make sure that we can say \emph{something} about how big $\abs{x+3}$ is.
So the first thing we'll do is to require that $\abs{x-3}<1$. Once we've
specified that $\abs{x-3}<1$, or $2<x<4$, we have $5<x+3<7$ and we've
guaranteed that $\abs{x+3}<7$. So we now have
\[
\abs{x^2-9}=\abs{x-3}\cdot\abs{x+3}<7\abs{x-3},
\]
which shows that we have $\abs{x^2-9}<\varepsilon$ for
$\abs{x-3}<\varepsilon/7$, provided that we've also required that
$\abs{x-3}<1$. Or, to make it look more official: we require that
$\abs{x-3}<\min(\varepsilon/7,1)$.

The initial specification $\abs{x-3}<1$ was simply made for convenience. We
could just as well have specified that $\abs{x-3}<\frac{1}{10}$ or
$\abs{x-3}<10$ or any other convenient number. To make sure you understand the
reasoning in the previous paragraph, it is a good exercise to figure out how
the argument would read if we chose $\abs{x-3}<10$.

Our argument to show that $f$ approaches $9$ near $3$ will basically work to
show that $f$ approaches $a^2$ near $a$ for any $a$, except that we need to
worry a bit more about getting the proper inequality for $\abs{x+a}$. We first
require that $\abs{x-a}<1$, again with the expectation that this will ensure
that $\abs{x+a}$ is not too large. In fact, Problem~1-12 shows that
\[
\abs{x}-\abs{a}\leq\abs{x-a}<1,
\]
so
\[
\abs{x}<1+\abs{a},
\]
and consequently
\[
\abs{x+a}\leq\abs{x}+\abs{a}<2\abs{a}+1,
\]
so that we then have
\begin{align*}
\abs{x^2-a^2}&=\abs{x-a}\cdot\abs{x+a}\\
&<\abs{x-a}\cdot(2\abs{a}+1),
\end{align*}
which shows that we have $\abs{x^2-a^2}<\varepsilon$ for
$\abs{x-a}<\varepsilon/(2\abs{a}+1)$, provided that we also have
$\abs{x-a}<1$. Officially: we require that
$\abs{x-a}<\min(\varepsilon/(2\abs{a}+1),1)$.

In contrast to this example, we'll now consider the function $f(x)=1/x$ (for
$x\neq 0$), and try to show that $f$ approaches $1/3$ near $3$. This means that
we need to show how to guarantee the inequality
\[
\abs*{\frac{1}{x}-\frac{1}{3}}<\varepsilon
\]
for any given positive number $\varepsilon$ by requiring $\abs{x-3}$ to be
small enough. We begin by writing
\[
\abs*{\frac{1}{x}-\frac{1}{3}}=\abs*{\frac{3-x}{3x}}=\frac{1}{3}\cdot\frac{1}{\abs{x}}\cdot\abs{x-3},
\]
giving us the nice factor $\abs{x-3}$, and even an extra $\frac{1}{3}$ for
good measure, along with the problem factor $1/\abs{x}$. In this case, we first
need to make sure that $\abs{x}$ isn't too \emph{small}, so that $1/\abs{x}$
won't be too large.

We can first require that $\abs{x-3}<1$, because this gives $2<x<4$, so that
\[
\frac{1}{4}<\frac{1}{x}<\frac{1}{2},
\]
which not only tells us that $\frac{1}{x}<\frac{1}{2}$, but also that $x>0$,
which is important in order to conclude that $\frac{1}{\abs{x}}<\frac{1}{2}$.
We now have
\[
\abs*{\frac{1}{x}-\frac{1}{3}}=\frac{1}{3}\cdot\frac{1}{\abs{x}}\cdot\abs{x-3}<\frac{1}{6}\abs{x-3},
\]
which shows that we have $\abs{1/x-1/3}<\varepsilon$ for
$\abs{x-3}<6\varepsilon$, provided that we've also required that
$\abs{x-3}<1$. Or, to make it look official again: we require that
$\abs{x-3}<\min(6\varepsilon,1)$.

If we instead wanted to show that $f$ approaches $-1/3$ near $-3$, we would
begin by stipulating that $\abs{x-(-3)}<1$, giving $-4<x<-2$, once again
implying that $\abs{1/x}<1/2$, so that everything works as before.

To show in general that $f$ approaches $1/a$ near $a$ for any $a$ we proceed
in basically the same way, except that, again, we have to be a little more
careful in formulating our initial stipulation. It's not good enough simply to
require that $\abs{x-a}$ should be less than $1$, or any other particular
number, because if $a$ is close to $0$ this would allow values of $x$ that are
negative (not to mention the embarrassing possibility that $x=0$, so that
$f(x)$ isn't even defined!).

The trick in this case is to first require that
\[
\abs{x-a}<\frac{\abs{a}}{2};
\]
in other words, we require that $x$ be less than half as far from $0$ as $a$
(Figure~7). You should be able to check first that $x\neq 0$ and that
$1/\abs{x}<2/\abs{a}$, and then work out the rest of the argument.

\begin{center}
\figtex{assets/geometric/ch05-fig07}\\[0.6em]
\figlabel{7}
\end{center}
With all the work required for these simple examples, you may have begun to
quail at the prospect of tackling even more complicated functions. But that
won't really be necessary, since we will eventually have some basic theorems
that we can rely on. Instead of worrying about the unpleasant algebra that
might be involved in functions like $f(x)=x^3$ or $f(x)=1/x^3$, we'll turn our
attention to some examples that might appear to be even more frightening.

Consider first the function $f(x)=x\sin 1/x$ (Figure~8). Despite the erratic
behavior of this function near $0$ it is clear, at least intuitively, that $f$
approaches $l=0$ near $a=0$ (remember that our provisional definition
specifically exempts $x=a$ from consideration, so it doesn't matter that this
function isn't even defined at $0$). We want to show that we can get
$f(x)=x\sin 1/x$ as close to $0$ as desired if we require that $x$ be
sufficiently close to $0$, but $\neq 0$. In other words, for any number
$\varepsilon>0$, we want to show that we can ensure that
\[
\abs{f(x)-0}=\abs*{x\sin\frac{1}{x}}<\varepsilon
\]
by requiring that $\abs{x}=\abs{x-0}$ is sufficiently small (but $\neq 0$).
But this is easy. Since
\[
\abs*{\sin\frac{1}{x}}\leq 1, \qquad\text{for all }x\neq 0,
\]
we have
\[
\abs*{x\sin\frac{1}{x}}\leq\abs{x}, \qquad\text{for all }x\neq 0,
\]
so we can make $\abs{x\sin 1/x}<\varepsilon$ simply by requiring that
$\abs{x}<\varepsilon$ and $\neq 0$.

\begin{center}
\figtex{assets/plots/ch05-fig08}\\[0.6em]
\figlabel{8}
\end{center}
For the function $f(x)=x^2\sin 1/x$ (Figure~9) it seems even clearer that $f$
approaches $0$ near $0$. If, for example, we want
\[
\abs*{x^2\sin\frac{1}{x}}<\frac{1}{10},
\]
then we certainly need only require that $\abs{x}<\frac{1}{10}$ and $x\neq 0$,
since this implies that $\abs{x^2}<\frac{1}{100}$ and consequently
\[
\abs*{x^2\sin\frac{1}{x}}\leq\abs{x^2}<\frac{1}{100}<\frac{1}{10}.
\]
(We could do even better, and allow $\abs{x}<1/\sqrt{10}$ and $x\neq 0$, but
there is no particular virtue in being as economical as possible.) In general,
if $\varepsilon>0$, to ensure that
\[
\abs*{x^2\sin\frac{1}{x}}<\varepsilon,
\]
we need only require that
\[
\abs{x}<\varepsilon \quad\text{and}\quad x\neq 0,
\]
provided that $\varepsilon\leq 1$. If we are given an $\varepsilon$ which is
greater than $1$ (it might be, even though it is ``small'' $\varepsilon$'s
which are of interest), then it does not suffice to require that
$\abs{x}<\varepsilon$, but it certainly suffices to require that
$\abs{x}<1$ and $x\neq 0$.

\begin{center}
\figtex{assets/plots/ch05-fig09}\\[0.6em]
\figlabel{9}
\end{center}
As a third example, consider the function $f(x)=\sqrt{\abs{x}}\sin 1/x$
(Figure~10). In order to make $\abs{\sqrt{\abs{x}}\sin 1/x}<\varepsilon$ we
can require that
\[
\abs{x}<\varepsilon^2 \quad\text{and}\quad x\neq 0
\]
(the algebra is left to you).

\begin{center}
\figtex{assets/plots/ch05-fig10}\\[0.6em]
\figlabel{10}
\end{center}
Finally, let us consider the function $f(x)=\sin 1/x$ (Figure~11). For this
function it is \emph{false} that $f$ approaches $0$ near $0$. This amounts to
saying that it is not true for every number $\varepsilon>0$ that we can get
$\abs{f(x)-0}<\varepsilon$ by choosing $x$ sufficiently small, and $\neq 0$.
To show this we simply have to find \emph{one} $\varepsilon>0$ for which the
condition $\abs{f(x)-0}<\varepsilon$ cannot be guaranteed, no matter how small
we require $\abs{x}$ to be. In fact, $\varepsilon=\frac{1}{2}$ will do: it is
impossible to ensure that $\abs{f(x)}<\frac{1}{2}$ no matter how small we
require $\abs{x}$ to be; for if $A$ is any interval containing $0$, there is
some number $x=1/(\frac{1}{2}\pi+2n\pi)$ which is in this interval, and for
this $x$ we have $f(x)=1$.

\begin{center}
\figtex{assets/plots/ch05-fig11}\\[0.6em]
\figlabel{11}
\end{center}
This same argument can be used (Figure~12) to show that $f$ does not approach
\emph{any} number near $0$. To show this we must again find, for any
particular number $l$, some number $\varepsilon>0$ so that
$\abs{f(x)-l}<\varepsilon$ is \emph{not} true, no matter how small $x$ is
required to be. The choice $\varepsilon=\frac{1}{2}$ works for any number $l$;
that is, no matter how small we require $\abs{x}$ to be, we cannot ensure that
$\abs{f(x)-l}<\frac{1}{2}$. The reason is, that for any interval $A$
containing $0$ we can find both $x_1$ and $x_2$ in this interval with
\[
f(x_1)=1 \qquad\text{and}\qquad f(x_2)=-1,
\]
namely
\[
x_1=\frac{1}{\frac{1}{2}\pi+2n\pi} \qquad\text{and}\qquad
x_2=\frac{1}{\frac{3}{2}\pi+2m\pi}
\]
for large enough $n$ and $m$. But the interval from $l-\frac{1}{2}$ to
$l+\frac{1}{2}$ cannot contain both $-1$ and $1$, since its total length is
only $1$; so we cannot have
\[
\abs{1-l}<\frac{1}{2} \quad\text{and also}\quad \abs{-1-l}<\frac{1}{2},
\]
no matter what $l$ is.

\begin{center}
\figtex{assets/plots/ch05-fig12}\\[0.6em]
\figlabel{12}
\end{center}
The phenomenon exhibited by $f(x)=\sin 1/x$ near $0$ can occur in many ways.
If we consider the function
\[
f(x)=\begin{cases}
0, & x\text{ irrational}\\
1, & x\text{ rational},
\end{cases}
\]
then, no matter what $a$ is, $f$ does not approach any number $l$ near $a$.
In fact, we cannot make $\abs{f(x)-l}<\frac{1}{4}$ no matter how close we
bring $x$ to $a$, because in any interval around $a$ there are numbers $x$
with $f(x)=0$, and also numbers $x$ with $f(x)=1$, so that we would need
$\abs{0-l}<\frac{1}{4}$ and also $\abs{1-l}<\frac{1}{4}$.

An amusing variation on this behavior is presented by the function shown in
Figure~13:
\[
f(x)=\begin{cases}
x, & x\text{ rational}\\
0, & x\text{ irrational}.
\end{cases}
\]
The behavior of this function is ``opposite'' to that of $g(x)=\sin 1/x$; it
approaches $0$ at $0$, but does not approach any number at $a$, if $a\neq 0$.
By now you should have no difficulty convincing yourself that this is true.

\begin{center}
\figtex{assets/geometric/ch05-fig13}\\[0.6em]
\figlabel{13}
\end{center}
We conclude with a very simple example (Figure~14):
\[
f(x)=\begin{cases}
-1, & x<0\\
1, & x>0.
\end{cases}
\]
If $a>0$, then $f$ approaches $1$ near $a$: indeed, to ensure that
$\abs{f(x)-1}<\varepsilon$ it certainly suffices to require that
$\abs{x-a}<a$, since this implies
\begin{align*}
&-a<x-a\\
\text{or}\quad &0<x
\end{align*}
so that $f(x)=1$. Similarly, if $b<0$, then $f$ approaches $-1$ near $b$: to
ensure that $\abs{f(x)-(-1)}<\varepsilon$ it suffices to require that
$\abs{x-b}<-b$. Finally, as you may easily check, $f$ does not approach any
number near $0$.

\begin{center}
\figtex{assets/geometric/ch05-fig14}\\[0.6em]
\figlabel{14}
\end{center}
The time has now come to point out that of the many demonstrations about
limits which we have given, not one has been a real proof. The fault lies not
with our reasoning, but with our definition. If our provisional definition of
a function was open to criticism, our provisional definition of approaching a
limit is even more vulnerable. This definition is simply not sufficiently
precise to be used in proofs. It is hardly clear how one ``makes'' $f(x)$
close to $l$ (whatever ``close'' means) by ``requiring'' $x$ to be
sufficiently close to $a$ (however close ``sufficiently'' close is supposed to
be). Despite the criticisms of our definition you may feel (I certainly hope
you do) that our arguments were nevertheless quite convincing. In order to
present any sort of argument at all, we have been practically forced to invent
the real definition. It is possible to arrive at this definition in several
steps, each one clarifying some obscure phrase which still remains. Let us
begin, once again, with the provisional definition:

\begin{quote}
The function $f$ approaches the limit $l$ near $a$, if we can make $f(x)$ as
close as we like to $l$ by requiring that $x$ be sufficiently close to, but
unequal to, $a$.
\end{quote}

The very first change which we made in this definition was to note that making
$f(x)$ close to $l$ meant making $\abs{f(x)-l}$ small, and similarly for $x$
and $a$:

\begin{quote}
The function $f$ approaches the limit $l$ near $a$, if we can make
$\abs{f(x)-l}$ as small as we like by requiring that $\abs{x-a}$ be
sufficiently small, and $x\neq a$.
\end{quote}

The second, more crucial, change was to note that making $\abs{f(x)-l}$ ``as
small as we like'' means making $\abs{f(x)-l}<\varepsilon$ for any
$\varepsilon>0$ that happens to be given us:

\begin{quote}
The function $f$ approaches the limit $l$ near $a$, if for every number
$\varepsilon>0$ we can make $\abs{f(x)-l}<\varepsilon$ by requiring that
$\abs{x-a}$ be sufficiently small, and $x\neq a$.
\end{quote}

There is a common pattern to all the demonstrations about limits which we have
given. For each number $\varepsilon>0$ we found some other positive number,
$\delta$ say, with the property that if $x\neq a$ and $\abs{x-a}<\delta$, then
$\abs{f(x)-l}<\varepsilon$. For the function $f(x)=x\sin 1/x$ (with $a=0$,
$l=0$), the number $\delta$ was just the number $\varepsilon$; for
$f(x)=\sqrt{\abs{x}}\sin 1/x$, it was $\varepsilon^2$; for $f(x)=x^2$ it was
the minimum of $1$ and $\varepsilon/(2\abs{a}+1)$. In general, it may not be
at all clear how to find the number $\delta$, given $\varepsilon$, but it is
the condition $\abs{x-a}<\delta$ which expresses how small ``sufficiently''
small must be:

\begin{quote}
The function $f$ approaches the limit $l$ near $a$, if for every
$\varepsilon>0$ there is some $\delta>0$ such that, for all $x$, if
$\abs{x-a}<\delta$ and $x\neq a$, then $\abs{f(x)-l}<\varepsilon$.
\end{quote}

This is practically the definition we will adopt. We will make only one
trivial change, noting that ``$\abs{x-a}<\delta$ and $x\neq a$'' can just as
well be expressed ``$0<\abs{x-a}<\delta$.''

\begin{spivakdefn}
The function $f$ \textbf{approaches the limit $l$ near $a$} means: for every
$\varepsilon>0$ there is some $\delta>0$ such that, for all $x$, if
$0<\abs{x-a}<\delta$, then $\abs{f(x)-l}<\varepsilon$.
\end{spivakdefn}

This definition is so important (\emph{everything} we do from now on depends
on it) that proceeding any further without knowing it is hopeless. If
necessary memorize it, like a poem! That, at least, is better than stating it
incorrectly; if you do this you are doomed to give incorrect proofs. A good
exercise in giving correct proofs is to review every fact already demonstrated
about functions approaching limits, giving formal proofs of each. In most
cases, this will merely involve a bit of rewording to make the arguments
conform to our formal definition---all the algebraic work has been done
already. When proving that $f$ does \emph{not} approach $l$ at $a$, be sure to
negate the definition correctly:

If it is \emph{not} true that
\begin{quote}
for every $\varepsilon>0$ there is some $\delta>0$ such that, for all $x$, if
$0<\abs{x-a}<\delta$, then $\abs{f(x)-l}<\varepsilon$,
\end{quote}
then
\begin{quote}
there is \emph{some} $\varepsilon>0$ such that for \emph{every} $\delta>0$
there is \emph{some} $x$ which satisfies $0<\abs{x-a}<\delta$ but not
$\abs{f(x)-l}<\varepsilon$.
\end{quote}

Thus, to show that the function $f(x)=\sin 1/x$ does not approach $0$ near
$0$, we consider $\varepsilon=\frac{1}{2}$ and note that for every
$\delta>0$ there is some $x$ with $0<\abs{x-0}<\delta$ but not
$\abs{\sin 1/x-0}<\frac{1}{2}$---namely, an $x$ of the form
$1/(\pi/2+2n\pi)$, where $n$ is so large that $1/(\pi/2+2n\pi)<\delta$.

As a final illustration of the use of the definition of a function approaching
a limit, we have reserved the function shown in Figure~15, a standard example,
but one of the most complicated:
\[
f(x)=\begin{cases}
0, & x\text{ irrational, }0<x<1\\
1/q, & x=p/q\text{ in lowest terms, }0<x<1.
\end{cases}
\]
(Recall that $p/q$ is in lowest terms if $p$ and $q$ are integers with no
common factor and $q>0$.)

\begin{center}
\figtex{assets/plots/ch05-fig15}\\[0.6em]
\figlabel{15}
\end{center}
For any number $a$, with $0<a<1$, the function $f$ approaches $0$ at $a$. To
prove this, consider any number $\varepsilon>0$. Let $n$ be a natural number
so large that $1/n\leq\varepsilon$. Notice that the only numbers $x$ for which
$\abs{f(x)-0}<\varepsilon$ could be false are:
\[
\frac{1}{2},\
\frac{1}{3},\frac{2}{3},\
\frac{1}{4},\frac{3}{4},\
\frac{1}{5},\frac{2}{5},\frac{3}{5},\frac{4}{5},\
\ldots,\
\frac{1}{n},\ldots,\frac{n-1}{n}.
\]
(If $a$ is rational, then $a$ might be one of these numbers.) However many of
these numbers there may be, there are, at any rate, only finitely many.
Therefore, of all these numbers, one is closest to $a$; that is,
$\abs{p/q-a}$ is smallest for one $p/q$ among these numbers. (If $a$ happens
to be one of these numbers, then consider only the values $\abs{p/q-a}$ for
$p/q\neq a$.) This closest distance may be chosen as the $\delta$. For if
$0<\abs{x-a}<\delta$, then $x$ is \emph{not} one of
\[
\frac{1}{2},\ldots,\frac{n-1}{n}
\]
and therefore $\abs{f(x)-0}<\varepsilon$ is true. This completes the proof.
Note that our description of the $\delta$ which works for a given
$\varepsilon$ is completely adequate---there is no reason why we must give a
formula for $\delta$ in terms of $\varepsilon$.

Armed with our definition, we are now prepared to prove our first theorem; you
have probably assumed the result all along, which is a very reasonable thing
to do. This theorem is really a test case for our definition: if the theorem
could not be proved, our definition would be useless.

\begin{theorem}
A function cannot approach two different limits near $a$. In other words, if
$f$ approaches $l$ near $a$, and $f$ approaches $m$ near $a$, then $l=m$.
\end{theorem}

\begin{proof}
Since this is our first theorem about limits it will certainly be necessary to
translate the hypotheses according to the definition.

Since $f$ approaches $l$ near $a$, we know that for any $\varepsilon>0$ there
is some number $\delta_1>0$ such that, for all $x$,
\[
\text{if }0<\abs{x-a}<\delta_1,\text{ then }\abs{f(x)-l}<\varepsilon.
\]
We also know, since $f$ approaches $m$ near $a$, that there is some
$\delta_2>0$ such that, for all $x$,
\[
\text{if }0<\abs{x-a}<\delta_2,\text{ then }\abs{f(x)-m}<\varepsilon.
\]
We have had to use two numbers, $\delta_1$ and $\delta_2$, since there is no
guarantee that the $\delta$ which works in one definition will work in the
other. But, in fact, it is now easy to conclude that for any $\varepsilon>0$
there is some $\delta>0$ such that, for all $x$,
\[
\text{if }0<\abs{x-a}<\delta,\text{ then }\abs{f(x)-l}<\varepsilon\text{ and }\abs{f(x)-m}<\varepsilon;
\]
we simply choose $\delta=\min(\delta_1,\delta_2)$.

To complete the proof we just have to pick a particular $\varepsilon>0$ for
which the two conditions
\[
\abs{f(x)-l}<\varepsilon \quad\text{and}\quad \abs{f(x)-m}<\varepsilon
\]
cannot both hold, if $l\neq m$. The proper choice is suggested by Figure~16.
If $l\neq m$, so that $\abs{l-m}>0$, we can choose $\abs{l-m}/2$ as our
$\varepsilon$. It follows that there is a $\delta>0$ such that, for all $x$,
\begin{align*}
\text{if }0<\abs{x-a}<\delta,\text{ then }&\abs{f(x)-l}<\frac{\abs{l-m}}{2}\\
\text{and }&\abs{f(x)-m}<\frac{\abs{l-m}}{2}.
\end{align*}
This implies that for $0<\abs{x-a}<\delta$ we have
\begin{align*}
\abs{l-m}&=\abs{l-f(x)+f(x)-m}\leq\abs{l-f(x)}+\abs{f(x)-m}\\
&<\frac{\abs{l-m}}{2}+\frac{\abs{l-m}}{2}\\
&=\abs{l-m},
\end{align*}
a contradiction.
\end{proof}

\begin{center}
\figtex{assets/geometric/ch05-fig16}\\[0.6em]
\figlabel{16}
\end{center}
The number $l$ which $f$ approaches near $a$ is denoted by
$\lim_{x\to a}f(x)$ (read: the limit of $f(x)$ as $x$ approaches $a$). This
definition is possible only because of Theorem~1, which ensures that
$\lim_{x\to a}f(x)$ never has to stand for two different numbers. The
equation
\[
\lim_{x\to a}f(x)=l
\]
has exactly the same meaning as the phrase
\[
\text{$f$ approaches $l$ near $a$.}
\]
The possibility still remains that $f$ does not approach $l$ near $a$, for any
$l$, so that $\lim_{x\to a}f(x)=l$ is false for every number $l$. This is
usually expressed by saying that ``$\lim_{x\to a}f(x)$ does not exist.''

Notice that our new notation introduces an extra, utterly irrelevant letter
$x$, which could be replaced by $t$, $y$, or any other letter which does not
already appear---the symbols
\[
\lim_{x\to a}f(x),\qquad \lim_{t\to a}f(t),\qquad \lim_{y\to a}f(y),
\]
all denote precisely the same number, which depends on $f$ and $a$, and has
nothing to do with $x$, $t$, or $y$ (these letters, in fact, do not denote
anything at all). A more logical symbol would be something like
$\operatorname*{lim}\limits_{a}f$, but this notation, despite its brevity, is
so infuriatingly rigid that almost no one has seriously tried to use it. The
notation $\lim_{x\to a}f(x)$ is much more useful because a function $f$ often
has no simple name, even though it might be possible to express $f(x)$ by a
simple formula involving $x$. Thus, the short symbol
\[
\lim_{x\to a}(x^2+\sin x)
\]
could be paraphrased only by the awkward expression
\[
\operatorname*{lim}\limits_{a}f,\text{ where }f(x)=x^2+\sin x.
\]
Another advantage of the standard symbolism is illustrated by the expressions
\begin{align*}
\lim_{x\to a}x+t^3,\\
\lim_{t\to a}x+t^3.
\end{align*}
The first means the number which $f$ approaches near $a$ when
\[
f(x)=x+t^3, \qquad\text{for all }x;
\]
the second means the number which $f$ approaches near $a$ when
\[
f(t)=x+t^3, \qquad\text{for all }t.
\]
You should have little difficulty (especially if you consult Theorem~2)
proving that
\begin{align*}
\lim_{x\to a}x+t^3&=a+t^3,\\
\lim_{t\to a}x+t^3&=x+a^3.
\end{align*}

These examples illustrate the main advantage of our notation, which is its
flexibility. In fact, the notation $\lim_{x\to a}f(x)$ is so flexible that
there is some danger of forgetting what it really means. Here is a simple
exercise in the use of this notation, which will be important later: first
interpret precisely, and then prove the equality of the expressions
\[
\lim_{x\to a}f(x) \qquad\text{and}\qquad \lim_{h\to 0}f(a+h).
\]

An important part of this chapter is the proof of a theorem which will make it
easy to find many limits, as we promised long ago. The proof depends upon
certain properties of inequalities and absolute values, hardly surprising when
one considers the definition of limit. Although these facts have already been
stated in Problems~1-20, 1-21, and 1-22, because of their importance they will
be presented once again, in the form of a lemma (a lemma is an auxiliary
theorem, a result that justifies its existence only by virtue of its prominent
role in the proof of another theorem). The lemma says, roughly, that if $x$ is
close to $x_0$, and $y$ is close to $y_0$, then $x+y$ will be close to
$x_0+y_0$, and $xy$ will be close to $x_0 y_0$, and $1/y$ will be close to
$1/y_0$. This intuitive statement is much easier to remember than the precise
estimates of the lemma, and it is not unreasonable to read the proof of
Theorem~2 first, in order to see just how these estimates are used.

\begin{spivakdefn}[LEMMA]
(1) If
\[
\abs{x-x_0}<\frac{\varepsilon}{2}\quad\text{and}\quad
\abs{y-y_0}<\frac{\varepsilon}{2},
\]
then
\[
\abs{(x+y)-(x_0+y_0)}<\varepsilon.
\]
(2) If
\[
\abs{x-x_0}<\min\Bigl(1,\frac{\varepsilon}{2(\abs{y_0}+1)}\Bigr)
\quad\text{and}\quad
\abs{y-y_0}<\frac{\varepsilon}{2(\abs{x_0}+1)},
\]
then
\[
\abs{xy-x_0 y_0}<\varepsilon.
\]
(3) If $y_0\neq 0$ and
\[
\abs{y-y_0}<\min\Bigl(\frac{\abs{y_0}}{2},\frac{\varepsilon\abs{y_0}^2}{2}\Bigr),
\]
then $y\neq 0$ and
\[
\abs*{\frac{1}{y}-\frac{1}{y_0}}<\varepsilon.
\]
\end{spivakdefn}

\begin{proof}
(1)
\begin{align*}
\abs{(x+y)-(x_0+y_0)}&=\abs{(x-x_0)+(y-y_0)}\\
&\leq\abs{x-x_0}+\abs{y-y_0}<\frac{\varepsilon}{2}+\frac{\varepsilon}{2}=\varepsilon.
\end{align*}

(2) Since $\abs{x-x_0}<1$ we have
\[
\abs{x}-\abs{x_0}\leq\abs{x-x_0}<1,
\]
so that
\[
\abs{x}<1+\abs{x_0}.
\]
Thus
\begin{align*}
\abs{xy-x_0 y_0}&=\abs{x(y-y_0)+y_0(x-x_0)}\\
&\leq\abs{x}\cdot\abs{y-y_0}+\abs{y_0}\cdot\abs{x-x_0}\\
&<(1+\abs{x_0})\cdot\frac{\varepsilon}{2(\abs{x_0}+1)}+\abs{y_0}\cdot\frac{\varepsilon}{2(\abs{y_0}+1)}\\
&<\frac{\varepsilon}{2}+\frac{\varepsilon}{2}=\varepsilon.
\end{align*}

(3) We have
\[
\abs{y_0}-\abs{y}\leq\abs{y-y_0}<\frac{\abs{y_0}}{2},
\]
so $\abs{y}>\abs{y_0}/2$. In particular, $y\neq 0$, and
\[
\frac{1}{\abs{y}}<\frac{2}{\abs{y_0}}.
\]
Thus
\[
\abs*{\frac{1}{y}-\frac{1}{y_0}}=\frac{\abs{y_0-y}}{\abs{y}\cdot\abs{y_0}}<\frac{2}{\abs{y_0}}\cdot\frac{1}{\abs{y_0}}\cdot\frac{\varepsilon\abs{y_0}^2}{2}=\varepsilon.
\]
\end{proof}

\begin{theorem}
If $\lim_{x\to a}f(x)=l$ and $\lim_{x\to a}g(x)=m$, then
\begin{align*}
&(1)& \lim_{x\to a}(f+g)(x)&=l+m;\\
&(2)& \lim_{x\to a}(f\cdot g)(x)&=l\cdot m.
\end{align*}
Moreover, if $m\neq 0$, then
\[
(3)\qquad \lim_{x\to a}\Bigl(\frac{1}{g}\Bigr)(x)=\frac{1}{m}.
\]
\end{theorem}

\begin{proof}
The hypothesis means that for every $\varepsilon>0$ there are
$\delta_1,\delta_2>0$ such that, for all $x$,
\begin{align*}
&\text{if }0<\abs{x-a}<\delta_1,\text{ then }\abs{f(x)-l}<\varepsilon,\\
\text{and}\quad &\text{if }0<\abs{x-a}<\delta_2,\text{ then }\abs{g(x)-m}<\varepsilon.
\end{align*}
This means (since, after all, $\varepsilon/2$ is also a positive number) that
there are $\delta_1,\delta_2>0$ such that, for all $x$,
\begin{align*}
&\text{if }0<\abs{x-a}<\delta_1,\text{ then }\abs{f(x)-l}<\frac{\varepsilon}{2},\\
\text{and}\quad &\text{if }0<\abs{x-a}<\delta_2,\text{ then }\abs{g(x)-m}<\frac{\varepsilon}{2}.
\end{align*}
Now let $\delta=\min(\delta_1,\delta_2)$. If $0<\abs{x-a}<\delta$, then
$0<\abs{x-a}<\delta_1$ and $0<\abs{x-a}<\delta_2$ are both true, so both
\[
\abs{f(x)-l}<\frac{\varepsilon}{2} \quad\text{and}\quad
\abs{g(x)-m}<\frac{\varepsilon}{2}
\]
are true. But by part~(1) of the lemma this implies that
$\abs{(f+g)(x)-(l+m)}<\varepsilon$. This proves~(1).

To prove~(2) we proceed similarly, after consulting part~(2) of the lemma. If
$\varepsilon>0$ there are $\delta_1,\delta_2>0$ such that, for all $x$,
\begin{align*}
&\text{if }0<\abs{x-a}<\delta_1,\text{ then }
\abs{f(x)-l}<\min\Bigl(1,\frac{\varepsilon}{2(\abs{m}+1)}\Bigr),\\
\text{and}\quad &\text{if }0<\abs{x-a}<\delta_2,\text{ then }
\abs{g(x)-m}<\frac{\varepsilon}{2(\abs{l}+1)}.
\end{align*}
Again let $\delta=\min(\delta_1,\delta_2)$. If $0<\abs{x-a}<\delta$, then
\[
\abs{f(x)-l}<\min\Bigl(1,\frac{\varepsilon}{2(\abs{m}+1)}\Bigr)
\quad\text{and}\quad
\abs{g(x)-m}<\frac{\varepsilon}{2(\abs{l}+1)}.
\]
So, by the lemma, $\abs{(f\cdot g)(x)-l\cdot m}<\varepsilon$, and this
proves~(2).

Finally, if $\varepsilon>0$ there is a $\delta>0$ such that, for all $x$,
\[
\text{if }0<\abs{x-a}<\delta,\text{ then }
\abs{g(x)-m}<\min\Bigl(\frac{\abs{m}}{2},\frac{\varepsilon\abs{m}^2}{2}\Bigr).
\]
But according to part~(3) of the lemma this means, first, that $g(x)\neq 0$,
so $(1/g)(x)$ makes sense, and second that
\[
\abs*{\Bigl(\frac{1}{g}\Bigr)(x)-\frac{1}{m}}<\varepsilon.
\]
This proves~(3).
\end{proof}

Using Theorem~2 we can prove, trivially, such facts as
\[
\lim_{x\to a}\frac{x^3+7x^5}{x^2+1}=\frac{a^3+7a^5}{a^2+1},
\]
without going through the laborious process of finding a $\delta$, given an
$\varepsilon$. We must begin with
\begin{align*}
\lim_{x\to a}7&=7,\\
\lim_{x\to a}1&=1,\\
\lim_{x\to a}x&=a,
\end{align*}
but these are easy to prove directly. If we \emph{want} to find the $\delta$,
however, the proof of Theorem~2 amounts to a prescription for doing this.
Suppose, to take a simpler example, that we want to find a $\delta$ such that,
for all $x$,
\[
\text{if }0<\abs{x-a}<\delta,\text{ then }\abs{x^2+x-(a^2+a)}<\varepsilon.
\]
Consulting the proof of Theorem~2(1), we see that we must first find
$\delta_1$ and $\delta_2>0$ such that, for all $x$,
\begin{align*}
&\text{if }0<\abs{x-a}<\delta_1,\text{ then }\abs{x^2-a^2}<\frac{\varepsilon}{2}\\
\text{and}\quad &\text{if }0<\abs{x-a}<\delta_2,\text{ then }\abs{x-a}<\frac{\varepsilon}{2}.
\end{align*}
Since we have already given proofs that $\lim_{x\to a}x^2=a^2$ and
$\lim_{x\to a}x=a$, we know how to do this:
\begin{align*}
\delta_1&=\min\Bigl(1,\frac{\varepsilon/2}{2\abs{a}+1}\Bigr),\\
\delta_2&=\frac{\varepsilon}{2}.
\end{align*}
Thus we can take
\[
\delta=\min(\delta_1,\delta_2)=\min\Bigl(\min\Bigl(1,\frac{\varepsilon/2}{2\abs{a}+1}\Bigr),\frac{\varepsilon}{2}\Bigr).
\]
If $a\neq 0$, the same method can be used to find a $\delta>0$ such that, for
all $x$,
\[
\text{if }0<\abs{x-a}<\delta,\text{ then }
\abs*{\frac{1}{x^2}-\frac{1}{a^2}}<\varepsilon.
\]
The proof of Theorem~2(3) shows that the second condition will follow if we
find a $\delta>0$ such that, for all $x$,
\[
\text{if }0<\abs{x-a}<\delta,\text{ then }
\abs{x^2-a^2}<\min\Bigl(\frac{\abs{a}^2}{2},\frac{\varepsilon\abs{a}^4}{2}\Bigr).
\]
Thus we can take
\[
\delta=\min\Biggl(1,\frac{\min\bigl(\abs{a}^2/2,\ \varepsilon\abs{a}^4/2\bigr)}{2\abs{a}+1}\Biggr).
\]
Naturally, these complicated expressions for $\delta$ can be simplified
considerably, after they have been derived.

One technical detail in the proof of Theorem~2 deserves some discussion. In
order for $\lim_{x\to a}f(x)$ to be defined it is, as we know, not necessary
for $f$ to be defined at $a$, nor is it necessary for $f$ to be defined at all
points $x\neq a$. However, there must be some $\delta>0$ such that $f(x)$ is
defined for $x$ satisfying $0<\abs{x-a}<\delta$; otherwise the clause
\[
\text{``if }0<\abs{x-a}<\delta,\text{ then }\abs{f(x)-l}<\varepsilon\text{''}
\]
would make no sense at all, since the symbol $f(x)$ would make no sense for
some $x$'s. If $f$ and $g$ are two functions for which the definition makes
sense, it is easy to see that the same is true for $f+g$ and $f\cdot g$. But
this is not so clear for $1/g$, since $1/g$ is undefined for $x$ with
$g(x)=0$. However, this fact gets established in the proof of Theorem~2(3).

There are times when we would like to speak of the limit which $f$ approaches
at $a$, even though there is no $\delta>0$ such that $f(x)$ is defined for $x$
satisfying $0<\abs{x-a}<\delta$. For example, we want to distinguish the
behavior of the two functions shown in Figure~17, even though they are not
defined for numbers less than $a$. For the function of Figure~17(a) we write
\[
\lim_{x\to a^+}f(x)=l \qquad\text{or}\qquad \lim_{x\downarrow a}f(x)=l.
\]

\begin{center}
\figtex{assets/geometric/ch05-fig17}\\[0.6em]
\figlabel{17}
\end{center}
(The symbols on the left are read: the limit of $f(x)$ as $x$ approaches $a$
from above.) These ``limits from above'' are obviously closely related to
ordinary limits, and the definition is very similar:
$\lim_{x\to a^+}f(x)=l$ means that for every $\varepsilon>0$ there is a
$\delta>0$ such that, for all $x$,
\[
\text{if }0<x-a<\delta,\text{ then }\abs{f(x)-l}<\varepsilon.
\]
(The condition ``$0<x-a<\delta$'' is equivalent to
``$0<\abs{x-a}<\delta$ and $x>a$.'')

``Limits from below'' (Figure~18) are defined similarly:
$\lim_{x\to a^-}f(x)=l$ (or $\lim_{x\uparrow a}f(x)=l$) means that for every
$\varepsilon>0$ there is a $\delta>0$ such that, for all $x$,
\[
\text{if }0<a-x<\delta,\text{ then }\abs{f(x)-l}<\varepsilon.
\]

\begin{center}
\figtex{assets/geometric/ch05-fig18}\\[0.6em]
\figlabel{18}
\end{center}
It is quite possible to consider limits from above and below even if $f$ is
defined for numbers both greater and less than $a$. Thus, for the function $f$
of Figure~14, we have
\[
\lim_{x\to 0^+}f(x)=1 \qquad\text{and}\qquad \lim_{x\to 0^-}f(x)=-1.
\]
It is an easy exercise (Problem~29) to show that $\lim_{x\to a}f(x)$ exists if
and only if $\lim_{x\to a^+}f(x)$ and $\lim_{x\to a^-}f(x)$ both exist and are
equal.

Like the definitions of limits from above and below, which have been smuggled
into the text informally, there are other modifications of the limit concept
which will be found useful. In Chapter~4 it was claimed that if $x$ is large,
then $\sin 1/x$ is close to $0$. This assertion is usually written
\[
\lim_{x\to\infty}\sin 1/x=0.
\]
The symbol $\lim_{x\to\infty}f(x)$ is read ``the limit of $f(x)$ as $x$
approaches $\infty$,'' or ``as $x$ becomes infinite,'' and a limit of the form
$\lim_{x\to\infty}f(x)$ is often called a limit at infinity.

Figure~19 illustrates a general situation where $\lim_{x\to\infty}f(x)=l$.
Formally, $\lim_{x\to\infty}f(x)=l$ means that for every $\varepsilon>0$ there
is a number $N$ such that, for all $x$,
\[
\text{if }x>N,\text{ then }\abs{f(x)-l}<\varepsilon.
\]
The analogy with the definition of ordinary limits should be clear: whereas
the condition ``$0<\abs{x-a}<\delta$'' expresses the fact that $x$ is close to
$a$, the condition ``$x>N$'' expresses the fact that $x$ is large.

\begin{center}
\figtex{assets/plots/ch05-fig19}\\[0.6em]
\figlabel{19}
\end{center}
We have spent so little time on limits from above and below, and at infinity,
because the general philosophy behind the definitions should be clear if you
understand the definition of ordinary limits (which are by far the most
important). Many exercises on these definitions are provided in the Problems,
which also contain several other types of limits which are occasionally
useful.

\subsection*{PROBLEMS}

\pnum{1.}
Find the following limits. (These limits all follow, after some algebraic
manipulations, from the various parts of Theorem~2; be sure you know which
ones are used in each case, but don't bother listing them.)
\begin{romans}
\item $\displaystyle\lim_{x\to 1}\dfrac{x^2-1}{x+1}$.
\item $\displaystyle\lim_{x\to 2}\dfrac{x^3-8}{x-2}$.
\item $\displaystyle\lim_{x\to 3}\dfrac{x^3-8}{x-2}$.
\item $\displaystyle\lim_{x\to y}\dfrac{x^n-y^n}{x-y}$.
\item $\displaystyle\lim_{y\to x}\dfrac{x^n-y^n}{x-y}$.
\item $\displaystyle\lim_{h\to 0}\dfrac{\sqrt{a+h}-\sqrt{a}}{h}$.
\end{romans}

\pnum{2.}
Find the following limits.
\begin{romans}
\item $\displaystyle\lim_{x\to 1}\dfrac{1-\sqrt{x}}{1-x}$.
\item $\displaystyle\lim_{x\to 0}\dfrac{1-\sqrt{1-x^2}}{x}$.
\item $\displaystyle\lim_{x\to 0}\dfrac{1-\sqrt{1-x^2}}{x^2}$.
\end{romans}

\pnum{3.}
In each of the following cases, determine the limit $l$ for the given $a$, and
prove that it is the limit by showing how to find a $\delta$ such that
$\abs{f(x)-l}<\varepsilon$ for all $x$ satisfying $0<\abs{x-a}<\delta$.
\begin{romans}
\item $f(x)=x\bigl[3-\cos(x^2)\bigr]$, $a=0$.
\item $f(x)=x^2+5x-2$, $a=2$.
\item $f(x)=\dfrac{100}{x}$, $a=1$.
\item $f(x)=x^4$, arbitrary $a$.
\item $f(x)=x^4+\dfrac{1}{x}$, $a=1$.
\item $f(x)=\dfrac{x}{2-\sin^2 x}$, $a=0$.
\item $f(x)=\sqrt{\abs{x}}$, $a=0$.
\item $f(x)=\sqrt{x}$, $a=1$.
\end{romans}

\pnum{4.}
For each of the functions in Problem~4-17, decide for which numbers $a$ the
limit $\lim_{x\to a}f(x)$ exists.

\pnum{*5.}
\begin{letters}
\item Do the same for each of the functions in Problem~4-19.
\item Same problem, if we use infinite decimals ending in a string of $0$'s
instead of those ending in a string of $9$'s.
\end{letters}

\pnum{6.}
Suppose the functions $f$ and $g$ have the following property: for all
$\varepsilon>0$ and all $x$,
\begin{align*}
\text{if }0<\abs{x-2}<\sin^2\Bigl(\frac{\varepsilon^2}{9}\Bigr)+\varepsilon,
&\text{ then }\abs{f(x)-2}<\varepsilon,\\
\text{if }0<\abs{x-2}<\varepsilon^2,
&\text{ then }\abs{g(x)-4}<\varepsilon.
\end{align*}
For each $\varepsilon>0$ find a $\delta>0$ such that, for all $x$,
\begin{romans}
\item if $0<\abs{x-2}<\delta$, then $\abs{f(x)+g(x)-6}<\varepsilon$.
\item if $0<\abs{x-2}<\delta$, then $\abs{f(x)g(x)-8}<\varepsilon$.
\item if $0<\abs{x-2}<\delta$, then
$\abs*{\dfrac{1}{g(x)}-\dfrac{1}{4}}<\varepsilon$.
\item if $0<\abs{x-2}<\delta$, then
$\abs*{\dfrac{f(x)}{g(x)}-\dfrac{1}{2}}<\varepsilon$.
\end{romans}

\pnum{7.}
Give an example of a function $f$ for which the following assertion is
\emph{false}: If $\abs{f(x)-l}<\varepsilon$ when $0<\abs{x-a}<\delta$, then
$\abs{f(x)-l}<\varepsilon/2$ when $0<\abs{x-a}<\delta/2$.

\pnum{8.}
\begin{letters}
\item If $\lim_{x\to a}f(x)$ and $\lim_{x\to a}g(x)$ do not exist, can
$\lim_{x\to a}[f(x)+g(x)]$ exist? Can $\lim_{x\to a}f(x)g(x)$ exist?
\item If $\lim_{x\to a}f(x)$ exists and $\lim_{x\to a}[f(x)+g(x)]$ exists,
must $\lim_{x\to a}g(x)$ exist?
\item If $\lim_{x\to a}f(x)$ exists and $\lim_{x\to a}g(x)$ does not exist,
can $\lim_{x\to a}[f(x)+g(x)]$ exist?
\item If $\lim_{x\to a}f(x)$ exists and $\lim_{x\to a}f(x)g(x)$ exists, does
it follow that $\lim_{x\to a}g(x)$ exists?
\end{letters}

\pnum{9.}
Prove that $\lim_{x\to a}f(x)=\lim_{h\to 0}f(a+h)$. (This is mainly an
exercise in understanding what the terms mean.)

\pnum{10.}
\begin{letters}
\item Prove that $\lim_{x\to a}f(x)=l$ if and only if
$\lim_{x\to a}[f(x)-l]=0$. (First see why the assertion is obvious; then
provide a rigorous proof. In this chapter most problems which ask for proofs
should be treated in the same way.)
\item Prove that $\lim_{x\to 0}f(x)=\lim_{x\to a}f(x-a)$.
\item Prove that $\lim_{x\to 0}f(x)=\lim_{x\to 0}f(x^3)$.
\item Give an example where $\lim_{x\to 0}f(x^2)$ exists, but
$\lim_{x\to 0}f(x)$ does not.
\end{letters}

\pnum{11.}
Suppose there is a $\delta>0$ such that $f(x)=g(x)$ when
$0<\abs{x-a}<\delta$. Prove that $\lim_{x\to a}f(x)=\lim_{x\to a}g(x)$. In
other words, $\lim_{x\to a}f(x)$ depends only on the values of $f(x)$ for $x$
near $a$---this fact is often expressed by saying that limits are a ``local
property.'' (It will clearly help to use $\delta'$, or some other letter,
instead of $\delta$, in the definition of limits.)

\pnum{12.}
\begin{letters}
\item Suppose that $f(x)\leq g(x)$ for all $x$. Prove that
$\lim_{x\to a}f(x)\leq\lim_{x\to a}g(x)$, provided that these limits exist.
\item How can the hypotheses be weakened?
\item If $f(x)<g(x)$ for all $x$, does it necessarily follow that
$\lim_{x\to a}f(x)<\lim_{x\to a}g(x)$?
\end{letters}

\pnum{13.}
Suppose that $f(x)\leq g(x)\leq h(x)$ and that
$\lim_{x\to a}f(x)=\lim_{x\to a}h(x)$. Prove that $\lim_{x\to a}g(x)$ exists,
and that $\lim_{x\to a}g(x)=\lim_{x\to a}f(x)=\lim_{x\to a}h(x)$. (Draw a
picture!)

\pnum{*14.}
\begin{letters}
\item Prove that if $\lim_{x\to 0}f(x)/x=l$ and $b\neq 0$, then
$\lim_{x\to 0}f(bx)/x=bl$. Hint: Write $f(bx)/x=b[f(bx)/bx]$.
\item What happens if $b=0$?
\item Part~(a) enables us to find $\lim_{x\to 0}(\sin 2x)/x$ in terms of
$\lim_{x\to 0}(\sin x)/x$. Find this limit in another way.
\end{letters}

\pnum{15.}
Evaluate the following limits in terms of the number
$\alpha=\lim_{x\to 0}(\sin x)/x$.
\begin{romans}
\item $\displaystyle\lim_{x\to 0}\dfrac{\sin 3x}{x}$.
\item $\displaystyle\lim_{x\to 0}\dfrac{\sin ax}{\sin bx}$.
\item $\displaystyle\lim_{x\to 0}\dfrac{\sin^2 2x}{x}$.
\item $\displaystyle\lim_{x\to 0}\dfrac{\sin^2 2x}{x^2}$.
\item $\displaystyle\lim_{x\to 0}\dfrac{1-\cos x}{x^2}$.
\item $\displaystyle\lim_{x\to 0}\dfrac{\tan^2 x+2x}{x+x^2}$.
\item $\displaystyle\lim_{x\to 0}\dfrac{x\sin x}{1-\cos x}$.
\item $\displaystyle\lim_{h\to 0}\dfrac{\sin(x+h)-\sin x}{h}$.
\item $\displaystyle\lim_{x\to 1}\dfrac{\sin(x^2-1)}{x-1}$.
\item $\displaystyle\lim_{x\to 0}\dfrac{x^2(3+\sin x)}{(x+\sin x)^2}$.
\item $\displaystyle\lim_{x\to 1}(x^2-1)^3\sin\Bigl(\dfrac{1}{x-1}\Bigr)^3$.
\end{romans}

\pnum{16.}
\begin{letters}
\item Prove that if $\lim_{x\to a}f(x)=l$, then
$\lim_{x\to a}\abs{f}(x)=\abs{l}$.
\item Prove that if $\lim_{x\to a}f(x)=l$ and $\lim_{x\to a}g(x)=m$, then
$\lim_{x\to a}\max(f,g)(x)=\max(l,m)$ and similarly for $\min$.
\end{letters}

\pnum{17.}
\begin{letters}
\item Prove that $\lim_{x\to 0}1/x$ does not exist, i.e., show that
$\lim_{x\to 0}1/x=l$ is false for every number $l$.
\item Prove that $\lim_{x\to 1}1/(x-1)$ does not exist.
\end{letters}

\pnum{18.}
Prove that if $\lim_{x\to a}f(x)=l$, then there is a number $\delta>0$ and a
number $M$ such that $\abs{f(x)}<M$ if $0<\abs{x-a}<\delta$. (What does this
mean pictorially?) Hint: Why does it suffice to prove that
$l-1<f(x)<l+1$ for $0<\abs{x-a}<\delta$?

\pnum{19.}
Prove that if $f(x)=0$ for irrational $x$ and $f(x)=1$ for rational $x$, then
$\lim_{x\to a}f(x)$ does not exist for any $a$.

\pnum{*20.}
Prove that if $f(x)=x$ for rational $x$, and $f(x)=-x$ for irrational $x$, then
$\lim_{x\to a}f(x)$ does not exist if $a\neq 0$.

\pnum{21.}
\begin{letters}
\item Prove that if $\lim_{x\to 0}g(x)=0$, then
$\lim_{x\to 0}g(x)\sin 1/x=0$.
\item Generalize this fact as follows: If $\lim_{x\to 0}g(x)=0$ and
$\abs{h(x)}\leq M$ for all $x$, then $\lim_{x\to 0}g(x)h(x)=0$. (Naturally it
is unnecessary to do part~(a) if you succeed in doing part~(b); actually the
statement of part~(b) may make it easier than~(a)---that's one of the values
of generalization.)
\end{letters}

\pnum{22.}
Consider a function $f$ with the following property: if $g$ is any function
for which $\lim_{x\to 0}g(x)$ does not exist, then
$\lim_{x\to 0}[f(x)+g(x)]$ also does not exist. Prove that this happens if and
only if $\lim_{x\to 0}f(x)$ \emph{does} exist. Hint: This is actually very
easy: the assumption that $\lim_{x\to 0}f(x)$ does not exist leads to an
immediate contradiction if you consider the right $g$.

\pnum{**23.}
This problem is the analogue of Problem~22 when $f+g$ is replaced by
$f\cdot g$. In this case the situation is considerably more complex, and the
analysis requires several steps (those in search of an especially challenging
problem can attempt an independent solution).
\begin{letters}
\item Suppose that $\lim_{x\to 0}f(x)$ exists and is $\neq 0$. Prove that if
$\lim_{x\to 0}g(x)$ does not exist, then $\lim_{x\to 0}f(x)g(x)$ also does
not exist.
\item Prove the same result if $\lim_{x\to 0}\abs{f(x)}=\infty$. (The precise
definition of this sort of limit is given in Problem~37.)
\item Prove that if neither of these two conditions holds, then there is a
function $g$ such that $\lim_{x\to 0}g(x)$ does not exist, but
$\lim_{x\to 0}f(x)g(x)$ does exist.

Hint: Consider separately the following two cases: (1)~for some
$\varepsilon>0$ we have $\abs{f(x)}>\varepsilon$ for all sufficiently small
$x$. (2)~For every $\varepsilon>0$, there are arbitrarily small $x$ with
$\abs{f(x)}<\varepsilon$. In the second case, begin by choosing points $x_n$
with $\abs{x_n}<1/n$ and $\abs{f(x_n)}<1/n$.
\end{letters}

\pnum{*24.}
Suppose that $A_n$ is, for each natural number $n$, some \emph{finite} set of
numbers in $[0,1]$, and that $A_n$ and $A_m$ have no members in common if
$m\neq n$. Define $f$ as follows:
\[
f(x)=\begin{cases}
1/n, & x\text{ in }A_n\\
0, & x\text{ not in }A_n\text{ for any }n.
\end{cases}
\]
Prove that $\lim_{x\to a}f(x)=0$ for all $a$ in $[0,1]$.

\pnum{25.}
Explain why the following definitions of $\lim_{x\to a}f(x)=l$ are all
correct: For every $\delta>0$ there is an $\varepsilon>0$ such that, for all
$x$,
\begin{romans}
\item if $0<\abs{x-a}<\varepsilon$, then $\abs{f(x)-l}<\delta$.
\item if $0<\abs{x-a}<\varepsilon$, then $\abs{f(x)-l}\leq\delta$.
\item if $0<\abs{x-a}<\varepsilon$, then $\abs{f(x)-l}<5\delta$.
\item if $0<\abs{x-a}<\varepsilon/10$, then $\abs{f(x)-l}<\delta$.
\end{romans}

\pnum{*26.}
Give examples to show that the following definitions of
$\lim_{x\to a}f(x)=l$ are \emph{not} correct.
\begin{letters}
\item For all $\delta>0$ there is an $\varepsilon>0$ such that if
$0<\abs{x-a}<\delta$, then $\abs{f(x)-l}<\varepsilon$.
\item For all $\varepsilon>0$ there is a $\delta>0$ such that if
$\abs{f(x)-l}<\varepsilon$, then $0<\abs{x-a}<\delta$.
\end{letters}

\pnum{27.}
For each of the functions in Problem~4-17 indicate for which numbers $a$ the
one-sided limits $\lim_{x\to a^+}f(x)$ and $\lim_{x\to a^-}f(x)$ exist.

\pnum{*28.}
\begin{letters}
\item Do the same for each of the functions in Problem~4-19.
\item Also consider what happens if decimals ending in $0$'s are used instead
of decimals ending in $9$'s.
\end{letters}

\pnum{29.}
Prove that $\lim_{x\to a}f(x)$ exists if
$\lim_{x\to a^+}f(x)=\lim_{x\to a^-}f(x)$.

\pnum{30.}
Prove that
\begin{romans}
\item $\displaystyle\lim_{x\to 0^+}f(x)=\lim_{x\to 0^-}f(-x)$.
\item $\displaystyle\lim_{x\to 0}f(\abs{x})=\lim_{x\to 0^+}f(x)$.
\item $\displaystyle\lim_{x\to 0}f(x^2)=\lim_{x\to 0^+}f(x)$.
\end{romans}
(These equations, and others like them, are open to several interpretations.
They might mean only that the two limits are equal if they both exist; or that
if a certain one of the limits exists, the other also exists and is equal to
it; or that if either limit exists, then the other exists and is equal to it.
Decide for yourself which interpretations are suitable.)

\pnum{31.}
Suppose that $\lim_{x\to a^-}f(x)<\lim_{x\to a^+}f(x)$. (Draw a picture to
illustrate this assertion.) Prove that there is some $\delta>0$ such that
$f(x)<f(y)$ whenever $x<a<y$ and $\abs{x-a}<\delta$ and $\abs{y-a}<\delta$.
Is the converse true?

\pnum{32.}
Prove that
\[
\lim_{x\to\infty}\frac{a_n x^n+\cdots+a_0}{b_m x^m+\cdots+b_0}
\]
(with $a_n\neq 0$ and $b_m\neq 0$) exists if and only if $m\geq n$. What is
the limit when $m=n$? When $m>n$? Hint: the one easy limit is
$\lim_{x\to\infty}1/x^k=0$; do some algebra so that this is the only
information you need.

\pnum{33.}
Find the following limits.
\begin{romans}
\item $\displaystyle\lim_{x\to\infty}\dfrac{x+\sin^3 x}{5x+6}$.
\item $\displaystyle\lim_{x\to\infty}\dfrac{x\sin x}{x^2+5}$.
\item $\displaystyle\lim_{x\to\infty}\sqrt{x^2+x}-x$.
\item $\displaystyle\lim_{x\to\infty}\dfrac{x^2(1+\sin^2 x)}{(x+\sin x)^2}$.
\end{romans}

\pnum{34.}
Prove that $\lim_{x\to 0^+}f(1/x)=\lim_{x\to\infty}f(x)$.

\pnum{35.}
Find the following limits in terms of the number
$\alpha=\lim_{x\to 0}(\sin x)/x$.
\begin{romans}
\item $\displaystyle\lim_{x\to\infty}\dfrac{\sin x}{x}$.
\item $\displaystyle\lim_{x\to\infty}x\sin\dfrac{1}{x}$.
\end{romans}

\pnum{36.}
Define ``$\lim_{x\to-\infty}f(x)=l$.''
\begin{letters}
\item Find $\lim_{x\to-\infty}(a_n x^n+\cdots+a_0)/(b_m x^m+\cdots+b_0)$.
\item Prove that $\lim_{x\to\infty}f(x)=\lim_{x\to-\infty}f(-x)$.
\item Prove that $\lim_{x\to 0^-}f(1/x)=\lim_{x\to-\infty}f(x)$.
\end{letters}

\pnum{37.}
We define $\lim_{x\to a}f(x)=\infty$ to mean that for all $N$ there is a
$\delta>0$ such that, for all $x$, if $0<\abs{x-a}<\delta$, then $f(x)>N$.
(Draw an appropriate picture!) (Of course, we may still say that
$\lim_{x\to a}f(x)$ ``does not exist'' in the usual sense.)
\begin{letters}
\item Show that $\lim_{x\to 3}1/(x-3)^2=\infty$.
\item Prove that if $f(x)>\varepsilon>0$ for all $x$, and
$\lim_{x\to a}g(x)=0$, then
\[
\lim_{x\to a}f(x)/\abs{g(x)}=\infty.
\]
\end{letters}

\pnum{38.}
\begin{letters}
\item Define $\lim_{x\to a^+}f(x)=\infty$ and
$\lim_{x\to a^-}f(x)=\infty$. (Or at least convince yourself that you could
write down the definitions if you had the energy. How many other such symbols
can you define?)
\item Prove that $\lim_{x\to 0^+}1/x=\infty$.
\item Prove that $\lim_{x\to 0^+}f(x)=\infty$ if and only if
$\lim_{x\to\infty}f(1/x)=\infty$.
\end{letters}

\pnum{39.}
Find the following limits, when they exist.
\begin{romans}
\item $\displaystyle\lim_{x\to\infty}\dfrac{x^3+4x-7}{7x^2-x+1}$.
\item $\displaystyle\lim_{x\to\infty}x(1+\sin^2 x)$.
\item $\displaystyle\lim_{x\to\infty}x\sin^2 x$.
\item $\displaystyle\lim_{x\to\infty}x^2\sin\dfrac{1}{x}$.
\item $\displaystyle\lim_{x\to\infty}\sqrt{x^2+2x}-x$.
\item $\displaystyle\lim_{x\to\infty}x\bigl(\sqrt{x+2}-\sqrt{x}\bigr)$.
\item $\displaystyle\lim_{x\to\infty}\dfrac{\sqrt{\abs{x}}}{x}$.
\end{romans}

\pnum{40.}
\begin{letters}
\item Find the perimeter of a regular $n$-gon inscribed in a circle of radius
$r$. [Answer: $2rn\sin(\pi/n)$.]
\item What value does this perimeter approach as $n$ becomes very large?
\item What limit can you guess from this?
\end{letters}

\pnum{41.}
\begin{letters}
\item For $c>1$, show that $c^{1/n}=\sqrt[n]{c}$ approaches $1$ as $n$ becomes
very large. Hint: Show that for any $\varepsilon>0$ we cannot have
$c^{1/n}>1+\varepsilon$ for large $n$.
\item More generally, if $c>0$, then $c^{1/n}$ approaches $1$ as $n$ becomes
very large.
\end{letters}

\pnum{*42.}
After sending the manuscript for the first edition of this book off to the
printer, I thought of a much simpler way to prove that $\lim_{x\to a}x^2=a^2$
and $\lim_{x\to a}x^3=a^3$, without going through all the factoring tricks on
page~92. Suppose, for example, that we want to prove that
$\lim_{x\to a}x^2=a^2$, where $a>0$. Given $\varepsilon>0$, we simply let
$\delta$ be the minimum of $\sqrt{a^2+\varepsilon}-a$ and
$a-\sqrt{a^2-\varepsilon}$ (see Figure~19); then $\abs{x-a}<\delta$ implies
that $\sqrt{a^2-\varepsilon}<x<\sqrt{a^2+\varepsilon}$, so
$a^2-\varepsilon<x^2<a^2+\varepsilon$, or $\abs{x^2-a^2}<\varepsilon$. It is
fortunate that these pages had already been set, so that I couldn't make these
changes, because this ``proof'' is completely fallacious. Wherein lies the
fallacy?

\begin{center}
\figtex{assets/plots/ch05-fig20}\\[0.6em]
\figlabel{20}
\end{center}